\begin{table}[htbp]
\caption{TDE dataset folder hierarchy with path descriptions.}
\label{tab:folder_structure}
\begin{tabular}{lrr}
\toprule
Path & Depth & Contents \\
\midrule
dataset/ & 0 & Root dataset directory \\
├── raw/ & 1 & Raw CT data root \\
│   ├── <specimen\_name>/ & 2 & Specimen taxon folder (e.g. archaeopteryx\_london) \\
│   │   └── <specimen\_id>/ & 3 & Acquisition ID folder (e.g. BMNH-37001) \\
│   │       ├── org\_slices/ & 4 & JP2 CT slice directory (canonical ESRF name) \\
│   │       │   └── *.jp2 & 5 & JPEG2000 CT slices (uint16, ordered) \\
│   │       └── mesh/ & 4 & 3-D surface mesh files (OBJ/PLY) \\
├── projections/ & 1 & 2-D projection root \\
│   └── <specimen>\_\_<id>/ & 2 & Specimen projection folder (<name>\_\_<id>) \\
│       ├── *\_ortho\_00.png & 3 & Orthographic projection PNG \\
│       └── *\_rot\_000.png & 3 & Rotational projection PNG \\
└── synthetic/ & 1 & Synthetic dataset root \\
├── synthetic\_images/ & 2 & Deformed projection images \\
│   └── *\_\_comp*.png & 3 & Example compression image \\
├── deformation\_masks/ & 2 & Binary deformation masks \\
│   └── *\_mask.png & 3 & Example mask file \\
└── synthetic\_labels.csv & 2 & Per-image label CSV (11 columns) \\
\bottomrule
\end{tabular}
\end{table}
